\origpage{251}{GRAPH MINORS. XII}
\hspace*{12pt}(7) \textit{The topological feasibility condition is satisfied.}

\textit{Subproof.} Take a drawing of the directed graph $L$ in $\Sigma$ and replace each vertex $v$ by a small closed disc $D_v$. By applying a homeomorphism of $\Sigma$ which maps the residue of each edge of $L$ onto $h_i$ with the appropriate orientation, the images of the $D_v$'s under this homeomorphism satisfy the topological feasibility condition.

From (5), (6), (7), and (3.2) we deduce that for each $v \in V(L)$ there is a connected subgraph $\alpha(v)$ of $G$ with $Z \cap V(\alpha(v)) = Z_v$, such that $\alpha(u), \alpha(v)$ are mutually disjoint for all distinct $u, v$. For $1 \le i \le t$ define $\alpha(e_i) = h_i$. Then $\alpha$ satisfies conditions (i)--(iv) in the definition of ``controls an $L$ minor,'' and we claim that it also satisfies (v). For certainly $\theta \ge |V(L)|$; suppose that $(A, B) \in \mathscr{T}$, and $V(\alpha(v)) \subseteq V(A)$ for some $v \in V(L)$. By our assumption, there exist distinct $i, j$ such that $v$ is incident with $e_i$ and $e_j$, and so $Z_v$ meets both $\{a_i, b_i\}$ and $\{a_j, b_j\}$. Consequently there exist $c_i, c_j \in Z_v$ with $d(c_i, c_j) \ge \theta_1$, since $d(a_i, h_i), d(b_i, h_i) \le 2$ and $d(a_j, h_j), d(b_j, h_j) \le 2$, and $d(h_i, h_j) \ge \theta_1 + 4$. Since $c_i, c_j \in Z_v \subseteq V(\alpha(v))$ and $\alpha(v)$ is connected and $V(\alpha(v)) \subseteq V(A)$, we deduce that there is a path $Q$ of $\Gamma$ with ends $c_i, c_j$ and with $V(Q) \subseteq V(A)$. By (4.1), $(A, B)$ has order $\ge \frac{1}{2} d(c_i, c_j) \ge \frac{1}{2} \theta_1 \ge |V(L)|$. Hence (v) holds, and the proof is complete. $\square$

\stmt
(4.4) \textit{For any connected surface $\Sigma$ and any integer $p \ge 0$, there is an integer $\theta \ge 1$ such that the following is true. Let $\Gamma$ be a 2-cell drawing in $\Sigma$, with a respectful tangle $\mathscr{T}$ of order $\ge \theta$. Let $\Gamma$ be a subgraph of a graph $G$, let $q = \frac{1}{2} p(p - 1)$, and for $1 \le i \le q$, let $P_i$ be a path of $G$ with distinct ends $u_i, v_i \in V(\Gamma)$ and with no other vertex or edge in $\Gamma$. Let $P_1, \dots, P_q$ be mutually disjoint, and for all distinct $u, v \in \{u_1, v_1, \dots, u_q, v_q\}$ let $d(u, v) \ge \theta$, where $d$ is the metric of $\mathscr{T}$. Then the tangle in $G$ induced by $\mathscr{T}$ controls a $K_p$ minor of $G$.}

\textit{Proof.} We may assume that $p \ge 5$, for otherwise $K_p$ can be drawn in $\Sigma$ and the result follows from (4.3). Let us choose $\theta' \ge 2p$ to satisfy (3.2) for $\Sigma$, taking $t = p(p - 1)$ and $z = p(p - 1)$. We claim that $\theta = \theta' + 2$ satisfies (4.4). For let $\Gamma$, $\mathscr{T}$, $P_1, \dots, P_q$ be as in the theorem. We may renumber $P_1, \dots, P_q$ as $P_{ij}$ ($1 \le i < j \le p$); and let $P_{ij}$ have ends $u_{ij}, v_{ij}$ ($1 \le i < j \le p$). Let $Z = \{u_{ij}, v_{ij} \colon 1 \le i < j \le p\}$, and for $1 \le j \le p$ let
\begin{equation*}
Z_j = \{u_{ij} \colon 1 \le i < j\} \cup \{v_{ij} \colon j < i \le p\}.
\end{equation*}
Then $|Z| = p(p - 1)$, and $(Z_1, \dots, Z_p)$ is a partition of $Z$. For each $v \in Z$, let $\Delta_v \subseteq \Sigma$ be a closed disc with $v \in \text{bd}(\Delta_v)$ and $U(\Gamma) \cap \Delta_v = \{v\}$, chosen all mutually disjoint. Let $\Delta_v \subseteq \{v\} \cup r_v$, where $r_v$ is a region. Since $V(\Gamma) \cap \text{bd}(\Delta_v)$ is free for each $v$ (because $\Gamma$ is connected) and the topological feasibility condition is trivially true (because $|V(\Gamma) \cap \text{bd}(\Delta_v)| = 1$ for each $v$), and $d(r_u, r_v) \ge d(u, v) - 2 \ge \theta - 2 = \theta'$, it follows from (3.2) that there
